\section{HAETAE 오류주입 취약성 분석}
\label{sec:analysis}

본 장에서는 \code{HAETAE}의 전체 서명 경로에서 발생 가능한 오류를
분류하고, 응답 연산의 항 제거가 비밀정보 노출로 이어지는 조건을
분석한다. 먼저 공격자가 유도하는 오류 효과와 관찰 가능한 출력을
정의한 뒤, T1, T2 및 거부 판정 우회에 대한 대수적 관계를 제시한다.

\begin{figure*}[t]
\centering
\begin{tikzpicture}[
  node distance=5mm,
  block/.style={
    draw,
    rounded corners=1pt,
    minimum height=12mm,
    text width=20mm,
    align=center,
    font=\small
  },
  fault/.style={
    draw,
    rounded corners=1pt,
    fill=black!7,
    minimum height=12mm,
    text width=22mm,
    align=center,
    font=\small
  },
  flow/.style={-{Latex[length=2mm]},thick}
]
\node[block] (sample) {$(\yvec,b)$\\샘플링};
\node[block,right=of sample] (challenge) {챌린지 $c$\\유도};
\node[fault,right=of challenge] (t1) {$c\svec_1$ 계산\\T1: 곱 항 제거};
\node[fault,right=of t1] (t2) {$+\yvec$ 계산\\T2: nonce 제거};
\node[fault,right=of t2] (rb) {거부 판정\\RB: 판정 우회};
\node[block,right=of rb] (output) {인코딩 및\\서명 출력};

\draw[flow] (sample) -- (challenge);
\draw[flow] (challenge) -- (t1);
\draw[flow] (t1) -- (t2);
\draw[flow] (t2) -- (rb);
\draw[flow] (rb) -- (output);

\node[below=7mm of t1,align=center,font=\footnotesize] (t1effect)
  {$\zvec_{\mathrm{T1}}\approx\yvec$\\동일 nonce 차분 필요};
\node[below=7mm of t2,align=center,font=\footnotesize] (t2effect)
  {$\zvec_{\mathrm{T2}}=\pm c\svec_1$\\단일 오류 응답 역산};
\node[below=7mm of rb,align=center,font=\footnotesize] (rbeffect)
  {경계 위반 응답\\잠재적 분포 누설};

\draw[-{Latex[length=1.6mm]},dashed] (t1) -- (t1effect);
\draw[-{Latex[length=1.6mm]},dashed] (t2) -- (t2effect);
\draw[-{Latex[length=1.6mm]},dashed] (rb) -- (rbeffect);
\end{tikzpicture}
\caption{Overview of the analyzed fault points in the HAETAE signing procedure}
\label{fig:fault-overview}
\end{figure*}


\subsection{오류 모델 및 분석 기준}

본 논문에서는 일시적 오류가 서명 1회에 영향을 주는 상황을 가정한다.
오류 효과는 데이터 변조와 프로그램 흐름 변조의 두 가지 유형으로
구분한다. 데이터 변조는 중간 버퍼나 연산 결과가 비정상 값으로
변경되는 오류이며, 프로그램 흐름 변조는 명령어 또는 반복문의 일부가
생략되는 오류이다. 물리적으로는 두 현상이 함께 발생할 수 있으나,
각 오류가 만드는 정보 노출 관계를 구분하기 위해 의미론적 효과를
기준으로 분석한다.

오류의 대수적 가능성과 검증 범위를 구분하기 위해 Table
~\ref{tab:evidence-level}의 분석 기준을 사용한다. E0은 코드 또는
수식에 공격 가능한 관계가 존재하는 단계이며, E1은 해당 관계에서
비밀정보를 계산하는 복원식이 도출된 단계이다. E2는 전체 서명 코드의
소프트웨어 오류 에뮬레이션에서 목표 효과를 재현한 경우를 의미한다.
E3는 실제 물리 글리치로 목표 중간값을 관측한 단계이며, E4는 디버그
출력 없이 직렬화 서명만으로 공격을 수행한 단계이다. 마지막으로 E5는
공개키 검증이나 서명 위조를 통해 최종 보안 영향을 확인한 경우이다.
이 분류는 논문 완성도를 순차적으로 판정하는 척도가 아니다. E2는
의미론적 오류가 만드는 취약성을 검증하는 데 충분한 증거이며, E3--E5는
특정 물리 구현과 공격 인터페이스에 대한 추가 검증 범위를 나타낸다.

\begin{table}[H]
\centering
\caption{Evidence levels for fault injection analysis}
\label{tab:evidence-level}
\footnotesize
\begin{tabularx}{\columnwidth}{@{}c l X@{}}
\toprule
Level & Objective & Required evidence \\
\midrule
E0 & 구조 분석 & 코드 또는 수식의 오류 관계 \\
E1 & 악용 가능성 & 공개정보를 이용한 복원식 \\
E2 & SW 재현 & 전체 서명 코드의 의미론적 오류 \\
E3 & 물리 재현 & 글리치로 생성된 목표 중간값 \\
E4 & 출력 공격 & 직렬화 서명만을 이용한 복원 \\
E5 & 영향 검증 & 공개키 검증 또는 위조 \\
\bottomrule
\end{tabularx}
\end{table}

이러한 구분은 SW 에뮬레이션의 취약성 증거를 물리 성공률로 오해하는
것을 방지하는 동시에, 물리 실험이 수행되지 않은 지점의 분석 결과를
불필요하게 축소하지 않도록 한다. 물리 실험에서 내부 계측을 사용한
경우에도 목표 오류 효과의 인과 검증이라는 범위는 유지된다.

\subsection{서명 경로의 오류 지점}

\code{HAETAE} 서명 과정의 오류 지점은 Table~\ref{tab:fault-surface}과
같다. SEED, SIGNBIT, UNPACK 및 LSB는 응답 계산 이전의 상류
연산을 대상으로 하며, 선행연구에서 제시된 지점이다
~\cite{lee2026haetae}. T1과 T2는 식~\eqref{eq:intro-response}의
두 항을 대상으로 한다. RB는 응답 생성 이후의 거부 판정을 생략하는
오류이다.

\begin{table*}[t]
\centering
\caption{Fault injection points in the HAETAE signing procedure}
\label{tab:fault-surface}
\small
\begin{tabularx}{0.96\textwidth}{@{}l X X l@{}}
\toprule
Point & Fault effect & Potential information or impact & Required output \\
\midrule
SEED & 샘플링 seed 고정 또는 변조 & 예측 가능한 nonce & 결함 서명 \\
SIGNBIT & bimodal 부호 변조 & 정상--결함 차분 관계 & 서명 쌍 \\
UNPACK & 공개행렬 언패킹 변조 & 정상--결함 차분 관계 & 서명 쌍 \\
LSB & 챌린지 유도값 변조 & 서로 다른 챌린지의 차분 & 서명 쌍 \\
T1 & $c\svec$ 항 제거 & nonce 및 차분 기반 $\svec_1$ 관계 & 동일 nonce 응답 쌍 \\
T2 & $\yvec$ 항 제거 & 단일 오류 응답의 $\svec_1$ 관계 & 결함 응답과 $c$ \\
RB & 거부 판정 우회 & 경계 위반과 잠재적 분포 누설 & 다수의 서명 전사 \\
\bottomrule
\end{tabularx}
\end{table*}

상류 오류는 난수 생성이나 챌린지 유도 과정을 변조하여 정상 서명과 다른
관계를 형성한다. 반면 T1과 T2는 비밀 성분과 nonce가 결합되는 응답
연산 자체를 대상으로 한다. 두 오류는 동일한 응답식의 서로 다른 항을
제거하지만 공격에 필요한 출력 수와 실현 조건이 다르다. RB는 즉각적인
대수적 복원식보다는 거부 샘플링이 보장하는 출력 분포를 변화시킨다는
점에서 구분된다.

본 연구의 SW 에뮬레이션은 각 지점에서 목표 연산을 생략하거나 값을
변조한 뒤 전체 서명 경로의 출력을 관찰한다. 따라서 Table
~\ref{tab:fault-surface}은 물리 발생 빈도가 아니라, 해당 의미론적
오류가 발생했을 때 형성되는 공격 관계를 나타낸다. 물리 검증은 이
가운데 T1을 주 대상으로 수행하며, T2 글리치 탐색은 비교 실험으로
제시한다.

\subsection{T2 오류를 이용한 비밀 성분 복원}
\label{subsec:t2-analysis}

T2는 응답 계산에서 일회용 난수 $\yvec$의 덧셈이 제거되는 오류이다.
깨끗한 T2 오류가 발생하면 응답은 다음과 같이 변한다.
\begin{equation}
  \zvec_{\mathrm{T2}}=(-1)^b c\svec_1 .
  \label{eq:t2-response}
\end{equation}
\code{HAETAE}의 챌린지 $c$는 서명에 포함되는 공개정보이므로, 공격자는
$c$와 오류 응답 사이의 다항식 곱 관계를 이용할 수 있다. 공개
챌린지의 NTT 표현을 $\widehat{c}$라 하고, $j$번째 비밀 다항식에
대응하는 오류 응답을 $\zvec_{\mathrm{T2}}^{(j)}$라 하면
$\widehat{c}[k]\neq0$인 성분에서 다음 식이 성립한다.
\begin{equation}
 \widehat{\svec}_1^{(j)}[k]
 =
 \operatorname{NTT}\!\left(
 \zvec_{\mathrm{T2}}^{(j)}\right)[k]
 \widehat{c}[k]^{-1}\pmod q .
 \label{eq:t2-recovery}
\end{equation}

부호 $(-1)^b$는 두 후보를 비교하거나 응답의 공개 성분을 이용하여
판별할 수 있다. 또한 내부 고정소수점 배율 $L_n=2^{13}$에 의해
$c\svec_1$이 정수 배율로 표현되므로, 깨끗한 항 제거가 발생한 경우
반올림 과정에서도 복원에 필요한 관계가 유지된다.

식~\eqref{eq:t2-recovery}는 T2의 대수적 악용 가능성을 보이지만,
임의의 단일 물리 글리치가 곧바로 비밀키 전체를 복구한다는 의미는
아니다. 복원 대상은 전체 서명키 패키지가 아니라 비밀 성분
$\svec_1$이다. 또한 글리치가 벡터 전체의 $\yvec$ 항만 제거해야 하며,
해당 응답이 인코딩 경로를 통과해 공격자가 관찰할 수 있는 형태로
출력되어야 한다.

\subsection{T1 오류를 이용한 차분 복원}
\label{subsec:t1-analysis}

T1은 $c\svec_1$ 곱 결과가 응답에 반영되지 않는 오류이다. 이 경우
결함 응답은 다음과 같이 nonce에 가까운 값이 된다.
\begin{equation}
  \zvec_{\mathrm{T1}}\approx\yvec .
  \label{eq:t1-response}
\end{equation}
동일한 nonce로 생성된 정상 응답을 $\zvec_{\mathrm{N}}$이라 하면 정상
응답과 결함 응답의 차분은
\begin{equation}
  \zvec_{\mathrm{N}}-\zvec_{\mathrm{T1}}
  =(-1)^b c\svec_1
  \label{eq:t1-difference}
\end{equation}
이 된다. 따라서 식~\eqref{eq:t2-recovery}와 동일한 NTT 역산을
적용하여 $\svec_1$을 복원할 수 있다.

T1은 T2와 달리 정상 응답과 결함 응답이 동일한 nonce를 사용해야 한다.
물리 글리치가 곱셈 반복의 일부만 생략한 경우에는 하나의 오류 응답에서
특정 다항식 블록만 복원될 수 있다. 이때 서로 다른 블록에 영향을 준
결함 응답을 누적하면 전체 $\svec_1$을 구성할 수 있다. 다만 고정
nonce는 반복 실험을 위한 통제 조건이므로, 이를 일반적인 배포 환경의
공격 조건과 구분해야 한다.

\subsection{거부 판정 우회}

RB는 응답이 서명 경계 $B_0$ 또는 $B_1$을 만족하지 않더라도 거부
분기를 수행하지 않는 오류이다. 이러한 응답이 방출되면 거부 샘플링이
의도한 출력 분포가 훼손될 수 있다. 또한 검증 단계는 별도의 경계
$B_2$를 사용하므로, 일부 경계 위반 응답이 대수적 검증식을 만족할
가능성이 있다.

그러나 경계 위반 응답의 존재만으로 비밀키 복원이 입증되는 것은 아니다.
인코더가 비정상 범위의 값을 다시 거부할 수 있으며, 통계적 누설을
주장하려면 다수의 출력에서 키 후보의 순위나 정보량을 제시해야 한다.
따라서 본 논문에서는 RB를 직접적인 키 복구 공격이 아니라 분포 누설을
가능하게 하는 잠재적 취약성으로 분류한다.
